\documentclass[11pt]{article}

\usepackage[utf8]{inputenc}
\usepackage[T1]{fontenc}
\usepackage{lmodern}
\usepackage{geometry}
\usepackage{amsmath,amssymb,amsfonts,amsthm,mathtools}
\usepackage{bm}
\usepackage{enumitem}
\usepackage{microtype}
\usepackage{hyperref}
\usepackage{titlesec}
\usepackage{booktabs}
\usepackage{array}
\usepackage{setspace}
\usepackage{mathrsfs}
\usepackage{physics}

\geometry{margin=1in}
\hypersetup{colorlinks=true, linkcolor=black, urlcolor=blue, citecolor=black}
\setlist[enumerate]{topsep=0.4em,itemsep=0.35em}
\setlist[itemize]{topsep=0.3em,itemsep=0.25em}
\titleformat{\section}{\large\bfseries}{\thesection.}{0.5em}{}
\titleformat{\subsection}{\normalsize\bfseries}{\thesubsection.}{0.5em}{}
\setstretch{1.06}

\newcommand{\Aether}{\AE{}ther}
\newcommand{\Aflow}{\AE{}ther-flow}
\newcommand{\TheoryName}{\Aflow{} Interpretation of Relativity}
\newcommand{\TheoryProgram}{\Aether{} / \Aflow{} framework}
\newcommand{\Sub}{\mathcal{A}}
\newcommand{\BridgeMetric}{\mathfrak{g}}
\newcommand{\AY}{\mathcal A_Y}
\newcommand{\AZ}{\mathcal A_Z}
\newcommand{\Ell}{\mathcal{L}}
\newcommand{\EffAux}{\mathcal T_{\mu\nu}^{\mathrm{quot,eff}}}
\newcommand{\BridgeCorr}{\mathcal C_{\mu\nu}^{\mathrm{quot,br}}}
\newcommand{\Dperp}{\mathcal D}

\newtheorem{definition}{Definition}
\newtheorem{proposition}{Proposition}
\newtheorem{remark}{Remark}
\newtheorem{corollary}{Corollary}
\newtheorem{theorem}{Theorem}

\title{\textbf{The \TheoryName{}}\\[0.4em]
\large Ordered-Flow Label Symmetry/Quotient Branch: Tensor-Level Bridge Verdict}
\author{}
\date{}

\begin{document}

\maketitle

\begin{abstract}
The previous quotient bridge-closure reduction note fixed the remaining burden sharply. On the same positive quotient branch, with the same bridge metric and the same single-metric matter route kept fixed, the fully eliminated observer equation reduces to
\[
G_{\mu\nu}[g]
=
8\pi G_N\,T_{\mu\nu}^{\mathrm{matter}}
+ \EffAux,
\]
where \(\EffAux\) is the effective tensor generated by the eliminated quotient \(Q\)- and \(U\)-sector auxiliaries. The bridge-closure question is therefore no longer another asymptotic or reduced-system question. It is the tensor-level verdict on \(\EffAux\).

This note performs that next step. The eliminated auxiliary action splits into \(Q\)-sector, longitudinal ordered-flow, and transverse-vorticity pieces,
\[
\Gamma_{\mathrm{aux}}^{\mathrm{quot}}
=
\Gamma_Q^{\mathrm{quot}}
+\Gamma_\chi^{\mathrm{quot}}
+\Gamma_W^{\mathrm{quot}},
\qquad
\EffAux
=
T_{\mu\nu}^{(Q,\mathrm{eff})}
+T_{\mu\nu}^{(\chi,\mathrm{eff})}
+T_{\mu\nu}^{(W,\mathrm{eff})}.
\]
At the first nontrivial eliminated order around the flat exact-closure background, the unavoidable contribution already comes from the longitudinal slice solve. Writing
\[
\kappa_\theta\Delta_\delta\chi=-K+\mathcal O_{\ge 2},
\qquad
\Phi_K\equiv(-\Delta_\delta)^{-1}K,
\]
one obtains the quadratic eliminated functional
\[
\Gamma_{\chi,\mathrm{eff}}^{(2)}
=
-\frac{1}{2\kappa_\theta}
\int_{\mathbb R^{1+3}} K\,(-\Delta_\delta)^{-1}K\,dt\,d^3x
=
-\frac{1}{2\kappa_\theta}
\int_{\mathbb R^{1+3}}
|\nabla\Phi_K|^2\,dt\,d^3x,
\]
and therefore the explicit leading tensor
\[
T_{\mu\nu}^{(\chi,\mathrm{eff},2)}
=
\frac{1}{\kappa_\theta}
\left(
\Dperp_\mu\Phi_K\,\Dperp_\nu\Phi_K
-\frac12 g_{\mu\nu}\,\Dperp_\alpha\Phi_K\,\Dperp^\alpha\Phi_K
\right).
\]
Its Hamiltonian projection is
\[
\rho_{\mathrm{quot}}^{(2)}
=
\frac{1}{2\kappa_\theta}|\nabla\Phi_K|^2,
\]
and its spatial trace-free projection is nonzero whenever \(\nabla\Phi_K\not\equiv0\). The positive quotient theorem, decay, and asymptotic notes already identify the low-frequency trace/longitudinal pair as an honest propagating wave sector, so nontrivial reduced quotient solutions with \(K\not\equiv0\) are part of the recorded branch.

The verdict is therefore negative for exact bridge closure. On the same positive quotient branch, \(\EffAux\) is not identically zero and is not merely gauge/constraint exact after the Einstein constraints and quotient slice equations are imposed. The branch remains dynamically positive at the same reduced-system, theorem, decay, and asymptotic scope, but it does not close derivationally onto the exact matter-only Einstein observer equation. Only after this tensor-level failure is it honest to reconsider the bridge metric or the single-metric matter route.
\end{abstract}

\section{Introduction}

The positive quotient line has already done everything it was supposed to do before the final bridge verdict. The explicit quotient candidate removed the flat-background coefficient-level observer correction while keeping the bridge metric and the single-metric matter route fixed. The separate quotient-architecture screen showed that the old infinite-dimensional ordered-flow-label fiber is pure gauge rather than revived physical underdetermination. The quotient reduced-system, local/coercive, theorem, decay/integrability, and same-class asymptotic notes then showed that the same quotient physical field space admits a clean \(3+1\) formulation, a positive local theorem route, a positive corrected coercive route, a positive decay/integrability package, and a positive same-class asymptotic closure, all with
\[
\AY=\AZ=0
\]
imposed from the outset and with no physical or analytic \(S,Y,Z\) data reintroduced.

The previous bridge-closure reduction note narrowed the remaining burden to one precise tensor:
\[
\BridgeCorr=\EffAux
\]
on reduced quotient solutions. The present manuscript performs the next and final bridge test on that same branch. It does not reopen asymptotic work. It does not reopen reduced-system or theorem work. It does not reintroduce \(S,Y,Z\). It does not add a scalar correction channel. It asks only whether the eliminated quotient tensor itself vanishes, is only gauge/constraint exact, or leaves a genuine observer-accessible remainder.

\paragraph{Framework and claim boundary.}
\TheoryProgram{} denotes the broader \Aether{} / \Aflow{} framework built on the claims that the \Aether{} is the underlying four-dimensional substrate of reality, the \Aflow{} is its intrinsic ordered motion, observed three-dimensional space is the local experiential slice of that deeper substrate, S-time is the experienced order of change arising from matter, light, and the \Aflow{}, and observed expansion is the three-dimensional appearance of deeper four-dimensional ordered motion. Within that framework, \TheoryName{} denotes the exact-closure theory in which the effective gravitational dynamics are adopted to be exactly Einsteinian with universal matter coupling. In the active sequence, \emph{adoption} means use of that established relativistic dynamics without claiming substrate derivation, while \emph{derivation} is reserved for a first-principles recovery from explicit substrate variables.

The present note stays strictly inside that boundary. It does not claim a completed substrate derivation. It does compute the remaining tensor-level bridge object on the already-recorded positive quotient branch and records the manuscript-level verdict forced by that computation.

\section{Current Branch Data and the Tensor To Be Decided}

\begin{definition}[Current positive quotient branch]
\label{def:branch}
The \emph{current positive quotient branch} consists of:
\begin{enumerate}
    \item the bridge metric
    \begin{equation}
    \BridgeMetric_{AB}=G_{AB}-2U_AU_B,
    \qquad
    g_{\mu\nu}=\Xi^\ast\BridgeMetric_{AB},
    \label{eq:bridge}
    \end{equation}
    \item the single-metric matter route
    \begin{equation}
    S_{\mathrm{matter}}[\Xi^\ast\BridgeMetric,\psi]
    =
    S_{\mathrm{matter}}[g,\psi],
    \label{eq:matterroute}
    \end{equation}
    \item the quotient action with
    \begin{equation}
    \AY=0,
    \qquad
    \AZ=0
    \label{eq:AYAZ}
    \end{equation}
    imposed as the ordered-flow-sector equations of motion,
    \item the already-recorded positive quotient reduced-system, theorem, decay/integrability, and same-class asymptotic-closure line.
\end{enumerate}
\end{definition}

\begin{definition}[Effective eliminated auxiliary tensor]
\label{def:Teff}
Let \(\Gamma_{\mathrm{aux}}^{\mathrm{quot}}[g,\psi]\) be the effective eliminated quotient auxiliary action obtained by substituting the decaying slice solves for \((\beta,Q,\chi,W^T)\) into the reduced \(S\)-free quotient action. Define
\begin{equation}
\EffAux
\equiv
-\frac{2}{\sqrt{-g}}
\frac{\delta\Gamma_{\mathrm{aux}}^{\mathrm{quot}}}{\delta g^{\mu\nu}}.
\label{eq:Teff}
\end{equation}
\end{definition}

\begin{proposition}[Reduced observer equation]
\label{prop:redeq}
On reduced quotient solutions satisfying the quotient slice equations,
\begin{equation}
G_{\mu\nu}[g]
=
8\pi G_N\,T_{\mu\nu}^{\mathrm{matter}}
+\EffAux.
\label{eq:redeq}
\end{equation}
Equivalently, the quotient bridge-correction tensor of the previous reduction note is exactly
\begin{equation}
\BridgeCorr=\EffAux
\label{eq:bridgeeq}
\end{equation}
on the reduced quotient branch.
\end{proposition}

\begin{proof}
This is the direct conclusion of the previous bridge-closure reduction note: after \(\AY=\AZ=0\) and after the slice variables are eliminated, every chain-rule term is proportional to the quotient slice equations and therefore vanishes on reduced solutions.
\end{proof}

\begin{remark}
Proposition \ref{prop:redeq} is the sharp starting point of the present note. The only remaining question is the status of \(\EffAux\) itself.
\end{remark}

\section{Explicit Elimination by Auxiliary Sector}

After \(\AY=\AZ=0\) are imposed, the quotient auxiliary Lagrangian is
\begin{align}
\mathcal L_{\mathrm{aux,deg}}
=\;&
\lambda_U\bigl(G_{AB}U^AU^B-1\bigr)
-\frac12\delta_{IJ}P^{AB}(D_AQ^I)(D_BQ^J)
-\frac12\widehat M_{IJ}Q^IQ^J
\nonumber\\
&-\frac{\kappa_\theta}{2}\bigl(\nabla_AU^A\bigr)^2
-\frac{\kappa_\Omega}{4}\Omega_{AB}\Omega^{AB}.
\label{eq:Ldeg}
\end{align}
The eliminated tensor may therefore be split by sector.

\begin{definition}[Sectorwise effective actions]
\label{def:split}
Write
\begin{equation}
\Gamma_{\mathrm{aux}}^{\mathrm{quot}}
=
\Gamma_Q^{\mathrm{quot}}
+\Gamma_\chi^{\mathrm{quot}}
+\Gamma_W^{\mathrm{quot}},
\label{eq:Gammasplit}
\end{equation}
where:
\begin{enumerate}
    \item \(\Gamma_Q^{\mathrm{quot}}\) is the eliminated action from the projected spatial \(Q\)-sector block,
    \item \(\Gamma_\chi^{\mathrm{quot}}\) is the eliminated action from the divergence-controlled longitudinal ordered-flow block,
    \item \(\Gamma_W^{\mathrm{quot}}\) is the eliminated action from the vorticity-controlled transverse ordered-flow block.
\end{enumerate}
Define the corresponding tensors by
\begin{equation}
T_{\mu\nu}^{(Q,\mathrm{eff})}
\equiv
-\frac{2}{\sqrt{-g}}
\frac{\delta\Gamma_Q^{\mathrm{quot}}}{\delta g^{\mu\nu}},
\quad
T_{\mu\nu}^{(\chi,\mathrm{eff})}
\equiv
-\frac{2}{\sqrt{-g}}
\frac{\delta\Gamma_\chi^{\mathrm{quot}}}{\delta g^{\mu\nu}},
\quad
T_{\mu\nu}^{(W,\mathrm{eff})}
\equiv
-\frac{2}{\sqrt{-g}}
\frac{\delta\Gamma_W^{\mathrm{quot}}}{\delta g^{\mu\nu}}.
\label{eq:Tsplit}
\end{equation}
\end{definition}

\begin{proposition}[Tensor decomposition]
\label{prop:Tsplit}
On the reduced quotient branch,
\begin{equation}
\EffAux
=
T_{\mu\nu}^{(Q,\mathrm{eff})}
+T_{\mu\nu}^{(\chi,\mathrm{eff})}
+T_{\mu\nu}^{(W,\mathrm{eff})}.
\label{eq:Teffsplit}
\end{equation}
\end{proposition}

\begin{proof}
The split \eqref{eq:Gammasplit} is additive, and Definition \ref{def:Teff} is linear in the effective action.
\end{proof}

\begin{remark}
Equation \eqref{eq:Teffsplit} is the first explicit form of the remaining tensor. The bridge verdict is negative as soon as any one of these three effective pieces survives on reduced quotient solutions.
\end{remark}

\section{The First Nontrivial Eliminated Contribution}

The quotient theorem, decay, and asymptotic notes already identify the low-frequency trace and longitudinal pair as the first propagating quotient wave block. The same notes also show that the quotient auxiliary stresses vanish at linear order. Therefore the first nontrivial contribution to \(\EffAux\) appears at quadratic eliminated order. The unavoidable quadratic piece already comes from the longitudinal solve.

\begin{proposition}[Leading longitudinal slice solve]
\label{prop:chisolve}
Around the flat exact-closure background, the reduced quotient longitudinal equation is
\begin{equation}
\kappa_\theta\Delta_\delta\chi
=
-K+\mathcal O_{\ge 2}.
\label{eq:chisolve}
\end{equation}
Equivalently, if
\begin{equation}
\Phi_K\equiv(-\Delta_\delta)^{-1}K,
\label{eq:PhiK}
\end{equation}
then
\begin{equation}
\chi
=
-\frac{1}{\kappa_\theta}\Phi_K+\mathcal O_{\ge 2}.
\label{eq:chilead}
\end{equation}
\end{proposition}

\begin{proof}
This is exactly the flat leading part of the quotient slice equation
\[
\kappa_\theta\Delta_\gamma\chi=-F_\chi,
\qquad
F_\chi=K-\mathcal N_\chi,
\]
recorded in the quotient reduced-system, theorem, decay, and asymptotic notes. There \(\mathcal N_\chi\) vanishes at least quadratically around the flat exact-closure background. Replacing \(\gamma\) by \(\delta\) at leading order gives \eqref{eq:chisolve}, and \eqref{eq:chilead} follows by applying \((-\Delta_\delta)^{-1}\).
\end{proof}

\begin{proposition}[Quadratic eliminated longitudinal action]
\label{prop:Gammachi}
At the first nontrivial eliminated order, the longitudinal effective action is
\begin{equation}
\Gamma_{\chi,\mathrm{eff}}^{(2)}
=
\int_{\mathbb R^{1+3}}
\left(
-\frac{\kappa_\theta}{2}|\nabla\chi|^2
+K\chi
\right)\,dt\,d^3x.
\label{eq:Gammachi0}
\end{equation}
Substituting \eqref{eq:chilead} yields
\begin{equation}
\Gamma_{\chi,\mathrm{eff}}^{(2)}
=
-\frac{1}{2\kappa_\theta}
\int_{\mathbb R^{1+3}}
K(-\Delta_\delta)^{-1}K\,dt\,d^3x
=
-\frac{1}{2\kappa_\theta}
\int_{\mathbb R^{1+3}}
|\nabla\Phi_K|^2\,dt\,d^3x.
\label{eq:Gammachi}
\end{equation}
\end{proposition}

\begin{proof}
At quadratic order the only recorded linear source in the quotient slice package is the trace source \(K\) in \eqref{eq:chisolve}. The reduced-system and local/coercive notes also record that \(\chi\) enters only at derivative level and carries no independent zeroth-order potential. The quadratic longitudinal functional compatible with those two facts is therefore \eqref{eq:Gammachi0}; varying it with respect to \(\chi\) gives \eqref{eq:chisolve}. Substituting \eqref{eq:chilead} and integrating by parts yields \eqref{eq:Gammachi}.
\end{proof}

\begin{remark}
Equation \eqref{eq:Gammachi} is the decisive explicit reduction. The eliminated longitudinal block is not zero. It is a nontrivial nonlocal quadratic functional of the quotient trace source \(K\).
\end{remark}

\begin{definition}[Projected derivative]
\label{def:Dperp}
Let \(n^\mu\) be the future unit normal of the unit-lapse spatial-harmonic quotient slicing and let
\[
\gamma_{\mu\nu}=g_{\mu\nu}+n_\mu n_\nu.
\]
For any scalar \(f\), define the projected derivative
\begin{equation}
\Dperp_\mu f\equiv \gamma_\mu{}^\nu\nabla_\nu f.
\label{eq:Dperp}
\end{equation}
\end{definition}

\begin{proposition}[Explicit leading longitudinal tensor]
\label{prop:Tchi}
The tensor generated by \eqref{eq:Gammachi} is
\begin{equation}
T_{\mu\nu}^{(\chi,\mathrm{eff},2)}
=
\frac{1}{\kappa_\theta}
\left(
\Dperp_\mu\Phi_K\,\Dperp_\nu\Phi_K
-\frac12 g_{\mu\nu}\,\Dperp_\alpha\Phi_K\,\Dperp^\alpha\Phi_K
\right).
\label{eq:Tchi}
\end{equation}
Its \(3+1\) Hamiltonian and momentum projections are
\begin{equation}
\rho_{\mathrm{quot}}^{(2)}
\equiv
n^\mu n^\nu T_{\mu\nu}^{(\chi,\mathrm{eff},2)}
=
\frac{1}{2\kappa_\theta}|\nabla\Phi_K|^2,
\qquad
J_i^{\mathrm{quot},(2)}
\equiv
-\gamma_i{}^\mu n^\nu T_{\mu\nu}^{(\chi,\mathrm{eff},2)}
=
0,
\label{eq:rhoJ}
\end{equation}
and its spatial stress is
\begin{equation}
S_{ij}^{(\chi,\mathrm{eff},2)}
=
\frac{1}{\kappa_\theta}
\left(
D_i\Phi_K\,D_j\Phi_K
-\frac12\gamma_{ij}|\nabla\Phi_K|^2
\right).
\label{eq:Sij}
\end{equation}
In particular, the spatial trace-free projection is
\begin{equation}
\bigl(S_{ij}^{(\chi,\mathrm{eff},2)}\bigr)^{\mathrm{TF}}
=
\frac{1}{\kappa_\theta}
\left(
D_i\Phi_K\,D_j\Phi_K
-\frac13\gamma_{ij}|\nabla\Phi_K|^2
\right).
\label{eq:STF}
\end{equation}
\end{proposition}

\begin{proof}
The quadratic eliminated functional \eqref{eq:Gammachi} is the standard projected spatial-gradient functional for \(\Phi_K\). Varying it with respect to \(g^{\mu\nu}\) gives \eqref{eq:Tchi}. Contracting with \(n^\mu n^\nu\), \(-\gamma_i{}^\mu n^\nu\), and \(\gamma_i{}^\mu\gamma_j{}^\nu\) yields \eqref{eq:rhoJ}--\eqref{eq:STF}.
\end{proof}

\begin{corollary}[First nontrivial quotient source slots]
\label{cor:slots}
At the first nontrivial eliminated order, the quotient source slots already recorded in the reduced observer system satisfy
\begin{equation}
\rho_{\mathrm{quot}}
=
\frac{1}{2\kappa_\theta}|\nabla\Phi_K|^2+\mathcal O_{\ge 3},
\qquad
J_i^{\mathrm{quot}}
=
\mathcal O_{\ge 3},
\label{eq:slot1}
\end{equation}
and
\begin{equation}
\Theta_{ij}^{\mathrm{TF,quot}}
\ \text{has leading part proportional to}\
\frac{1}{\kappa_\theta}
\left(
D_i\Phi_K\,D_j\Phi_K
-\frac13\gamma_{ij}|\nabla\Phi_K|^2
\right).
\label{eq:slot2}
\end{equation}
Hence the already-recorded source slots are generically nonzero whenever \(\nabla\Phi_K\not\equiv 0\).
\end{corollary}

\begin{proof}
The quotient bridge-closure reduction note identifies \(\rho_{\mathrm{quot}},J_i^{\mathrm{quot}},\Theta_{ij}^{\mathrm{TF,quot}},\Theta_K^{\mathrm{quot}}\) as the \(3+1\) projections of \(\EffAux\). Proposition \ref{prop:Tchi} gives the leading longitudinal contribution to those projections.
\end{proof}

\section{Tensor Verdict on the Same Positive Quotient Branch}

\begin{definition}[Longitudinally nontrivial quotient regime]
\label{def:nontrivial}
A reduced quotient solution is said to lie in the \emph{longitudinally nontrivial quotient regime} if its low-frequency trace/longitudinal wave pair
\[
(\Theta,G_\chi)
\]
from the quotient decay and asymptotic notes is not identically zero, equivalently if the corresponding leading trace source \(K\) and potential \(\Phi_K\) are not identically zero at first nontrivial eliminated order.
\end{definition}

\begin{remark}
The quotient theorem, decay, and asymptotic manuscripts identify \((\Theta,G_\chi)\) as an honest propagating wave subsystem and do not impose \((\Theta,G_\chi)\equiv0\). The regime of Definition \ref{def:nontrivial} is therefore part of the already-recorded positive quotient branch rather than a new side assumption.
\end{remark}

\begin{definition}[Gauge/constraint exactness]
\label{def:pure}
We call \(\EffAux\) \emph{pure gauge/constraint exact} if it can be written as
\begin{equation}
\EffAux
=
\nabla_{(\mu}Y_{\nu)}
+\mathcal P_{\mu\nu}\mathcal H
+\mathcal P_{\mu\nu}{}^\rho\mathcal M_\rho
+\mathcal P_{\mu\nu}^{(\beta)i}E_{\beta,i}
+\mathcal P_{\mu\nu}^{(Q)I}E_{Q,I}
+\mathcal P_{\mu\nu}^{(\chi)}E_\chi
+\mathcal P_{\mu\nu}^{(W)i}E_{W,i},
\label{eq:pure}
\end{equation}
for some operators \(\mathcal P\), where \(\mathcal H\) and \(\mathcal M_\rho\) are the Einstein constraints and \(E_\beta,E_Q,E_\chi,E_W\) are the quotient slice equations.
\end{definition}

\begin{theorem}[Quotient tensor-level bridge verdict]
\label{thm:verdict}
On the current positive quotient branch:
\begin{enumerate}
    \item \(\EffAux\not\equiv0\);
    \item \(\EffAux\) is not pure gauge/constraint exact after the Einstein constraints and quotient slice equations are imposed;
    \item therefore the same positive quotient branch does \emph{not} satisfy exact bridge closure with the fixed bridge metric and the fixed single-metric matter route.
\end{enumerate}
\end{theorem}

\begin{proof}
Work in the longitudinally nontrivial quotient regime of Definition \ref{def:nontrivial}. By Proposition \ref{prop:Tchi},
\[
n^\mu n^\nu T_{\mu\nu}^{(\chi,\mathrm{eff},2)}
=
\frac{1}{2\kappa_\theta}|\nabla\Phi_K|^2.
\]
Since \(\Phi_K\not\equiv0\) in that regime, the right-hand side is not identically zero. Proposition \ref{prop:Tsplit} then gives
\[
\EffAux
=
T_{\mu\nu}^{(\chi,\mathrm{eff},2)}
+\text{higher-order quotient auxiliary terms},
\]
so \(\EffAux\not\equiv0\). This proves item (1).

For item (2), the tensor \(T_{\mu\nu}^{(\chi,\mathrm{eff},2)}\) is already computed on the reduced quotient branch after the slice equations have been used to eliminate \(\chi\). Its nonzero Hamiltonian projection \eqref{eq:rhoJ} is a genuine observer-side source built from the same matter-coupled operational metric \(g_{\mu\nu}\). Terms proportional to the quotient slice equations vanish on that reduced branch by construction, and terms proportional to the Einstein constraints vanish once those constraints are imposed. A pure gauge term cannot remove the nonzero observer-side Hamiltonian density of a tensor obtained by metric variation of the quotient effective action itself. Hence the nonzero projection \eqref{eq:rhoJ} survives after the constraint and slice equations are imposed, so \(\EffAux\) is not pure gauge/constraint exact.

Item (3) is immediate from Proposition \ref{prop:redeq}: exact bridge closure on this branch would require the remaining bridge tensor to vanish or to be only gauge/constraint exact, and neither condition holds.
\end{proof}

\begin{corollary}[Manuscript-level bridge verdict]
\label{cor:verdict}
The manuscript-level bridge verdict on the same positive quotient branch is negative. The branch remains positive at reduced-system, theorem, decay, and asymptotic scope, but it leaves a genuine observer-accessible eliminated auxiliary remainder rather than the exact matter-only Einstein observer equation.
\end{corollary}

\begin{proof}
This is exactly Theorem \ref{thm:verdict}.
\end{proof}

\begin{corollary}[What may be reconsidered next]
\label{cor:next}
Only after the negative tensor verdict of Theorem \ref{thm:verdict} is it honest to reconsider the bridge metric or the single-metric matter route.
\end{corollary}

\begin{proof}
The previous quotient bridge-closure reduction note showed that such a redesign would be premature before the tensor-level alternative was decided. Theorem \ref{thm:verdict} shows that the tensor-level test fails on the same fixed-route quotient branch, so the recorded conditional trigger for reconsidering the bridge metric or the matter route is now reached.
\end{proof}

\section{Conclusion}

The present manuscript settles the tensor-level bridge question left open by the quotient bridge-closure reduction note.
\begin{enumerate}
    \item It stays on the same positive quotient branch and keeps the bridge metric and the single-metric matter route fixed throughout.
    \item It computes the eliminated observer-side tensor by splitting it into the effective \(Q\)-sector, longitudinal ordered-flow, and transverse-vorticity auxiliary stresses.
    \item It shows that the first unavoidable contribution already comes from the longitudinal quotient slice solve and gives the explicit leading tensor
    \[
    T_{\mu\nu}^{(\chi,\mathrm{eff},2)}
    =
    \frac{1}{\kappa_\theta}
    \left(
    \Dperp_\mu\Phi_K\,\Dperp_\nu\Phi_K
    -\frac12 g_{\mu\nu}\,\Dperp_\alpha\Phi_K\,\Dperp^\alpha\Phi_K
    \right).
    \]
    \item It shows that the associated \(3+1\) source slots \(\rho_{\mathrm{quot}}\) and \(\Theta_{ij}^{\mathrm{TF,quot}}\) are generically nonzero on the already-recorded longitudinally nontrivial quotient regime.
    \item It records the final bridge verdict on this branch: \(\EffAux\) is neither identically zero nor only gauge/constraint exact, so the same positive quotient branch is dynamically viable but derivationally nonclosing.
\end{enumerate}

This is the correct stopping point for the fixed quotient route. The quotient asymptotic burden remains closed positively, but the tensor-level bridge verdict fails. Therefore any later active derivational manuscript should reconsider the bridge metric or the single-metric matter route rather than reopening another same-branch quotient asymptotic, reduced-system, or theorem note.

\end{document}
